\begin{textsample}{POS Dim 1 – 1994 – Score 15.00 – 1994/t002458.txt}  \label{ex:f1_pos_001}
Because I usually take \textbf{the} role of trying to explain to people how wonderful \textbf{the} new technologies that are coming along are going to be, and I thought that, since I was among friends here, I would tell you what I really think and try to look back and try to understand what is really going on here with these amazing jumps in technology that seem so fast that we can barely keep on top of it.
So I ' m going to start out by showing just one very boring technology \textbf{slide}.
And then, so if you can just turn on \textbf{the} \textbf{slide} that ' s on.
This is just a random \textbf{slide} that I picked out of my file.
What I want to show you is not so much \textbf{the} details of \textbf{the} \textbf{slide}, but \textbf{the} general form of it.
This happens to be a \textbf{slide} of some analysis that we were doing about \textbf{the} power of RISC microprocessors versus \textbf{the} power of local area networks.
And \textbf{the} interesting thing about it is that this \textbf{slide}, like so many technology \textbf{slides} that we ' re used to, is a sort of a straight line on a semi-log curve.
In other words, every step here represents an order of magnitude in performance scale.
And this is a new thing that we talk about technology on semi-log curves.
Something really weird is going on here.
And that ' s basically what I ' m going to be talking about.
So, if you could bring up \textbf{the} lights.
If you could bring up \textbf{the} lights higher, because I ' m just going to use a piece of paper here.
Now why do we draw technology curves in semi-log curves?
Well \textbf{the} answer is, if I drew it on a normal curve where, let ' s say, this is years, this is time of some sort, and this is whatever measure of \textbf{the} technology that I ' m trying to graph, \textbf{the} graphs look sort of silly.
They sort of go like this.
And they don ' t tell us much.
Now if I graph, for instance, some other technology, say \textbf{transportation} technology, on a semi-log curve, it would look very stupid, it would look like a flat line.
But when something like this happens, things are qualitatively changing.
So if \textbf{transportation} technology was moving along as fast as microprocessor technology, then \textbf{the} day after tomorrow, I would be able to get in a taxi cab and be in Tokyo in 30 seconds.
It ' s not moving like that.
And there ' s nothing precedented in \textbf{the} history of technology development of this kind of self- \textbf{feeding} growth where you go by orders of magnitude every few years.
Now \textbf{the} question that I ' d like to ask is, if you look at these exponential curves, they don ' t go on forever.
Things just can ' t possibly keep changing as fast as they are.
One of two things is going to happen.
Either it ' s going to turn into a sort of classical S-curve like this, until something totally different comes along, or maybe it ' s going to do this.
That ' s about all it can do.
Now I ' m an optimist, so I sort of think it ' s probably going to do something like that.
If so, that means that what we ' re in \textbf{the} middle of right now is a transition.
We ' re sort of on this line in a transition from \textbf{the} way \textbf{the} world used to be to some new way that \textbf{the} world is.
And so what I ' m trying to ask, what I ' ve been asking myself, is what ' s this new way that \textbf{the} world is?
What ' s that new state that \textbf{the} world is heading toward?
Because \textbf{the} transition seems very, very confusing when we ' re right in \textbf{the} middle of it.
Now when I was a kid growing up, \textbf{the} future was kind of \textbf{the} year 2000, and people used to talk about what would happen in \textbf{the} year 2000.
Now here ' s a conference in which people talk about \textbf{the} future, and you notice that \textbf{the} future is still at about \textbf{the} year 2000.
It ' s about as far as we go out.
So in other words, \textbf{the} future has kind of been shrinking one year per year for my whole lifetime.
Now I think that \textbf{the} reason is because we all feel that something ' s happening there.
That transition is happening.
We can all sense it.
And we know that it just doesn ' t make too much sense to think out 30, 50 years because everything ' s going to be so different that a simple extrapolation of what we ' re doing just doesn ' t make any sense at all.
So what I would like to talk about is what that could be, what that transition could be that we ' re going through.
Now in order to do that I ' m going to have to talk about a bunch of stuff that really has nothing to do with technology and \textbf{computers}.
Because I think \textbf{the} only way to understand this is to really step back and take a long time scale look at things.
So \textbf{the} time scale that I would like to look at this on is \textbf{the} time scale of life on Earth.
So I think this picture makes sense if you look at it a few billion years at a time.
So if you go back about two and a half billion years, \textbf{the} Earth was this big, sterile hunk of \textbf{rock} with a lot of chemicals floating around on it.
And if you look at \textbf{the} way that \textbf{the} chemicals got organized, we begin to get a pretty good idea of how they do it.
And I think that there ' s theories that are beginning to understand about how it started with RNA, but I ' m going to tell a sort of simple story of it, which is that, at that time, there were little drops of \textbf{oil} floating around with all kinds of different recipes of chemicals in them.
And some of those drops of \textbf{oil} had a particular combination of chemicals in them which caused them to incorporate chemicals from \textbf{the} outside and grow \textbf{the} drops of \textbf{oil}.
And those that were like that started to split and divide.
And those were \textbf{the} most primitive forms of cells in a sense, those little drops of \textbf{oil}.
But now those drops of \textbf{oil} weren ' t really alive, as we say it now, because every one of them was a little random recipe of chemicals.
And every time it divided, they got sort of unequal division of \textbf{the} chemicals within them.
And so every drop was a little bit different.
In fact, \textbf{the} drops that were different in a way that caused them to be better at incorporating chemicals around them, grew more and incorporated more chemicals and divided more.
So those tended to live longer, get expressed more.
Now that ' s sort of just a very simple chemical form of life, but when things got interesting was when these drops learned a trick about abstraction.
Somehow by ways that we don ' t quite understand, these little drops learned to write down information.
They learned to record \textbf{the} information that was \textbf{the} recipe of \textbf{the} cell onto a particular kind of chemical called DNA.
So in other words, they worked out, in this mindless sort of \textbf{evolutionary} way, a form of writing that let them write down what they were, so that that way of writing it down could get \textbf{copied}. \textbf{The} amazing thing is that that way of writing seems to have stayed steady since it \textbf{evolved} two and a half billion years ago.
In fact \textbf{the} recipe for us, our genes, is exactly that same code and that same way of writing.
In fact, every living creature is written in exactly \textbf{the} same set of letters and \textbf{the} same code.
In fact, one of \textbf{the} things that I did just for amusement purposes is we can now write things in this code.
And I ' ve got here a little 100 micrograms of white \textbf{powder}, which I try not to let \textbf{the} security people see at airports.
But this has in it — what I did is I took this code — \textbf{the} code has standard letters that we use for symbolizing it — and I wrote my business card onto a piece of DNA and amplified it 10 to \textbf{the} 22 times.
So if anyone would like a hundred million \textbf{copies} of my business card, I have plenty for everyone in \textbf{the} room, and, in fact, everyone in \textbf{the} world, and it ' s right here.
If I had really been a egotist, I would have put it into a virus and released it in \textbf{the} room.
So what was \textbf{the} next step?
Writing down \textbf{the} DNA was an interesting step.
And that caused these cells — that kept them happy for another billion years.
But then there was another really interesting step where things became completely different, which is these cells started exchanging and communicating information, so that they began to get communities of cells.
I don ' t know if you know this, but bacteria can actually exchange DNA.
Now that ' s why, for instance, antibiotic resistance has \textbf{evolved}.
Some bacteria figured out how to stay away from penicillin, and it went around sort of creating its little DNA information with other bacteria, and now we have a lot of bacteria that are resistant to penicillin, because bacteria communicate.
Now what this communication allowed was communities to form that, in some sense, were in \textbf{the} same boat together ; they were synergistic.
So they survived or they failed together, which means that if a community was very successful, all \textbf{the} individuals in that community were repeated more and they were favored by \textbf{evolution}.
Now \textbf{the} transition point happened when these communities got so close that, in fact, they got together and decided to write down \textbf{the} whole recipe for \textbf{the} community together on one string of DNA.
And so \textbf{the} next stage that ' s interesting in life took about another billion years.
And at that stage, we have multi-cellular communities, communities of lots of different types of cells, working together as a single organism.
And in fact, we ' re such a multi- cellular community.
We have lots of cells that are not out for themselves anymore.
Your skin cell is really useless without a heart cell, muscle cell, a brain cell and so on.
So these communities began to \textbf{evolve} so that \textbf{the} interesting level on which \textbf{evolution} was taking place was no longer a cell, but a community which we call an organism.
Now \textbf{the} next step that happened is within these communities.
These communities of cells, again, began to \textbf{abstract} information.
And they began building very special structures that did nothing but process information within \textbf{the} community.
And those are \textbf{the} neural structures.
So neurons are \textbf{the} information processing apparatus that those communities of cells built up.
And in fact, they began to get specialists in \textbf{the} community and special structures that were responsible for recording, understanding, learning information.
And that was \textbf{the} brains and \textbf{the} nervous system of those communities.
And that gave them an \textbf{evolutionary} advantage.
Because at that point, an individual — \textbf{learning} could happen within \textbf{the} time span of a single organism, instead of over this \textbf{evolutionary} time span.
So an organism could, for instance, learn not to eat a certain kind of fruit because it tasted bad and it got sick last time it ate it.
That could happen within \textbf{the} lifetime of a single organism, whereas before they ' d built these special information processing structures, that would have had to be learned evolutionarily over hundreds of thousands of years by \textbf{the} individuals dying off that ate that kind of fruit.
So that nervous system, \textbf{the} fact that they built these special information structures, tremendously sped up \textbf{the} whole process of \textbf{evolution}.
Because \textbf{evolution} could now happen within an individual.
It could happen in learning time scales.
But then what happened was \textbf{the} individuals worked out, of course, tricks of communicating.
And for example, \textbf{the} most sophisticated version that we ' re aware of is human language.
It ' s really a pretty amazing invention if you think about it.
Here I have a very complicated, messy, confused idea in my head.
I ' m sitting here making grunting sounds basically, and hopefully constructing a similar messy, confused idea in your head that bears some analogy to it.
But we ' re taking something very complicated, turning it into sound, sequences of sounds, and producing something very complicated in your brain.
So this allows us now to begin to start functioning as a single organism.
And so, in fact, what we ' ve done is we, humanity, have started \textbf{abstracting} out.
We ' re going through \textbf{the} same levels that multi-cellular organisms have gone through — \textbf{abstracting} out our methods of recording, presenting, processing information.
So for example, \textbf{the} invention of language was a tiny step in that direction.
Telephony, \textbf{computers}, videotapes, CD-ROMs and so on are all our specialized mechanisms that we ' ve now built within our society for handling that information.
And it all connects us together into something that is much bigger and much faster and able to \textbf{evolve} than what we were before.
So now, \textbf{evolution} can take place on a scale of microseconds.
And you saw Ty ' s little \textbf{evolutionary} example where he sort of did a little bit of \textbf{evolution} on \textbf{the} Convolution program right before your eyes.
So now we ' ve speeded up \textbf{the} time scales once again.
So \textbf{the} first steps of \textbf{the} story that I told you about took a billion years a piece.
And \textbf{the} next steps, like nervous systems and brains, took a few hundred million years.
Then \textbf{the} next steps, like language and so on, took less than a million years.
And these next steps, like electronics, seem to be taking only a few decades. \textbf{The} process is feeding on itself and becoming, I guess, autocatalytic is \textbf{the} word for it — when something reinforces its rate of change. \textbf{The} more it changes, \textbf{the} faster it changes.
And I think that that ' s what we ' re seeing here in this \textbf{explosion} of curve.
We ' re seeing this process feeding back on itself.
Now I design \textbf{computers} for a living, and I know that \textbf{the} mechanisms that I use to design \textbf{computers} would be impossible without recent advances in \textbf{computers}.
So right now, what I do is I design \textbf{objects} at such complexity that it ' s really impossible for me to design them in \textbf{the} traditional sense.
I don ' t know what every transistor in \textbf{the} connection machine does.
There are billions of them.
Instead, what I do and what \textbf{the} designers at Thinking Machines do is we think at some level of abstraction and then we hand it to \textbf{the} machine and \textbf{the} machine takes it beyond what we could ever do, much farther and faster than we could ever do.
And in fact, sometimes it takes it by methods that we don ' t quite even understand.
One method that ' s particularly interesting that I ' ve been using a lot lately is \textbf{evolution} itself.
So what we do is we put inside \textbf{the} machine a process of \textbf{evolution} that takes place on \textbf{the} microsecond time scale.
So for example, in \textbf{the} most extreme cases, we can actually \textbf{evolve} a program by starting out with random sequences of instructions.
Say,"Computer, would you please make a hundred million random sequences of instructions.
Now would you please run all of those random sequences of instructions, run all of those programs, and pick out \textbf{the} ones that came closest to doing what I wanted."So in other words, I define what I wanted.
Let ' s say I want to sort numbers, as a simple example I ' ve done it with.
So find \textbf{the} programs that come closest to sorting numbers.
So of course, random sequences of instructions are very unlikely to sort numbers, so none of them will really do it.
But one of them, by luck, may put two numbers in \textbf{the} right order.
And I say,"Computer, would you please now take \textbf{the} 10 percent of those random sequences that did \textbf{the} best job.
Save those.
Kill off \textbf{the} rest.
And now let ' s reproduce \textbf{the} ones that sorted numbers \textbf{the} best.
And let ' s reproduce them by a process of recombination analogous to sex."Take two programs and they produce children by exchanging their subroutines, and \textbf{the} children inherit \textbf{the} traits of \textbf{the} subroutines of \textbf{the} two programs.
So I ' ve got now a new generation of programs that are produced by combinations of \textbf{the} programs that did a little bit better job.
Say,"Please repeat that process."Score them again.
Introduce some mutations perhaps.
And try that again and do that for another generation.
Well every one of those generations just takes a few milliseconds.
So I can do \textbf{the} equivalent of millions of years of \textbf{evolution} on that within \textbf{the} \textbf{computer} in a few minutes, or in \textbf{the} complicated cases, in a few hours.
At \textbf{the} end of that, I end up with programs that are absolutely perfect at sorting numbers.
In fact, they are programs that are much more efficient than programs I could have ever written by hand.
Now if I look at those programs, I can ' t tell you how they work.
I ' ve tried looking at them and telling you how they work.
They ' re \textbf{obscure}, weird programs.
But they do \textbf{the} job.
And in fact, I know, I ' m very confident that they do \textbf{the} job because they come from a line of hundreds of thousands of programs that did \textbf{the} job.
In fact, their life depended on doing \textbf{the} job.
I was riding in a 747 with Marvin Minsky once, and he pulls out this card and says,"Oh look.
Look at this.
It says, ' This plane has hundreds of thousands of tiny parts working together to make you a safe flight. ' Doesn ' t that make you feel confident?"In fact, we know that \textbf{the} engineering process doesn ' t work very well when it gets complicated.
So we ' re beginning to depend on \textbf{computers} to do a process that ' s very different than engineering.
And it lets us produce things of much more complexity than normal engineering lets us produce.
And yet, we don ' t quite understand \textbf{the} options of it.
So in a sense, it ' s getting ahead of us.
We ' re now using those programs to make much faster \textbf{computers} so that we ' ll be able to run this process much faster.
So it ' s \textbf{feeding} back on itself. \textbf{The} thing is becoming faster and that ' s why I think it seems so confusing.
Because all of these technologies are feeding back on themselves.
We ' re taking off.
And what we are is we ' re at a point in time which is analogous to when single- celled organisms were turning into multi-celled organisms.
So we ' re \textbf{the} amoebas and we can ' t quite figure out what \textbf{the} hell this thing is we ' re creating.
We ' re right at that point of transition.
But I think that there really is something coming along after us.
I think it ' s very haughty of us to think that we ' re \textbf{the} end product of \textbf{evolution}.
And I think all of us here are a part of producing whatever that next thing is.
So lunch is coming along, and I think I will stop at that point, before I get selected out.

% matched lemmas: abstract, computer, copy, evolution, evolutionary, evolve, explosion, feeding, learning, object, obscure, oil, powder, rock, slide, the, transportation
\end{textsample}
